\section{End-Term Examination --- Set A}
\label{sec:endterm_set_a}

\LPUPaperHeader{End-Term Examination}{A}{3 Hours}{60 Marks}{CO1 to CO6 (Units 1 to 6 Comprehensive)}

\begin{center}
\sffamily\textbf{\color{AlertRed}Structure of Question Paper:}\\[1mm]
\textbf{Part A (Objective):} 30 MCQs strictly from \textbf{Units 4, 5, and 6}. Total weight: \textbf{20 Marks}. All questions compulsory.\\
\textbf{Part B (Subjective):} 6 Questions from \textbf{all units} (10 Marks each). \textbf{Attempt any 4 questions} ($4 \times 10 = \mathbf{40\text{ Marks}}$).\\
\textbf{Total Maximum Marks: 60 Marks.}
\end{center}

\vspace{3mm}

% ------------------------------------------------------------------------------
% PART A: 30 MCQS FROM UNITS 4, 5, 6
% ------------------------------------------------------------------------------
\subsection*{PART A: Objective Section (30 MCQs --- Units 4, 5, 6) [20 Marks]}

\mcqitem{1}{The value of $\lim_{(x,y)\to(0,0)} \frac{xy}{\sqrt{x^2 + y^2}}$ is:}{$0$}{$1$}{Does not exist}{$\infty$}{1}

\mcqitem{2}{If $u = \ln(x^3 + y^3 + z^3 - 3xyz)$, then $\frac{\partial u}{\partial x} + \frac{\partial u}{\partial y} + \frac{\partial u}{\partial z}$ is:}{$\frac{3}{x+y+z}$}{$\frac{1}{x+y+z}$}{$\frac{9}{x+y+z}$}{$0$}{1}

\mcqitem{3}{If $f(x,y)$ is homogeneous of degree $n$, then $x^2 f_{xx} + 2xy f_{xy} + y^2 f_{yy}$ equals:}{$n(n-1)f$}{$n^2 f$}{$n f$}{$(n+1)f$}{1}

\mcqitem{4}{The Jacobian $\frac{\partial(x,y)}{\partial(r,\theta)}$ for polar coordinates $x = r\cos\theta, y = r\sin\theta$ is:}{$r$}{$r^2$}{$1/r$}{$1$}{1}

\mcqitem{5}{At a critical point $(a,b)$, if $rt - s^2 < 0$, the point is a:}{Local Minimum}{Local Maximum}{Saddle Point}{Point of Inflexion}{1}

\mcqitem{6}{The directional derivative of $\phi = xy + yz + zx$ at $(1, 2, 0)$ in the direction of $\hat{i} + 2\hat{j} + 2\hat{k}$ is:}{$\frac{10}{3}$}{$3$}{$\frac{8}{3}$}{$4$}{1}

\mcqitem{7}{A vector field $\vec{F}$ is called \textbf{solenoidal} if:}{$\nabla \times \vec{F} = 0$}{$\nabla \cdot \vec{F} = 0$}{$\nabla^2 \vec{F} = 0$}{$\vec{F} \cdot \vec{r} = 0$}{1}

\mcqitem{8}{A vector field $\vec{F}$ is called \textbf{irrotational} if:}{$\nabla \cdot \vec{F} = 0$}{$\nabla \times \vec{F} = \vec{0}$}{$\nabla \phi = \vec{F}$}{$\nabla^2 \vec{F} = 0$}{1}

\mcqitem{9}{The value of $\nabla \cdot \vec{r}$ where $\vec{r} = x\hat{i} + y\hat{j} + z\hat{k}$ is:}{$0$}{$1$}{$2$}{$3$}{1}

\mcqitem{10}{The value of $\nabla \times \vec{r}$ is:}{$\vec{0}$}{$3$}{$\vec{r}$}{$1$}{1}

\mcqitem{11}{The value of the double integral $\int_0^1 \int_0^2 xy \, dx \, dy$ is:}{$1$}{$2$}{$1/2$}{$4$}{1}

\mcqitem{12}{The area of a region $R$ in the $xy$-plane is represented by:}{$\iint_R dx \, dy$}{$\iint_R xy \, dx \, dy$}{$\iint_R (x+y) \, dx \, dy$}{$\oint_C x \, dx$}{1}

\mcqitem{13}{In spherical polar coordinates, the volume element $dV$ is:}{$\rho^2 \sin\phi \, d\rho \, d\phi \, d\theta$}{$\rho \sin\phi \, d\rho \, d\phi \, d\theta$}{$\rho^2 \cos\phi \, d\rho \, d\phi \, d\theta$}{$\rho^2 \, d\rho \, d\phi \, d\theta$}{1}

\mcqitem{14}{Green's Theorem in a plane converts a line integral around a closed curve into a:}{Surface Integral}{Double Integral over the enclosed region}{Triple Integral}{Volume Integral}{1}

\mcqitem{15}{Stokes' Theorem relates a line integral around a boundary curve $C$ to a:}{Volume Integral}{Surface Integral of $\nabla \times \vec{F}$}{Double Integral of $\nabla \cdot \vec{F}$}{Triple Integral}{1}

\mcqitem{16}{Gauss Divergence Theorem relates a surface integral over a closed surface $S$ to a:}{Line Integral}{Volume Integral of $\nabla \cdot \vec{F}$}{Surface Integral of Curl}{Double Integral}{1}

\mcqitem{17}{The value of $\nabla(r^n)$ where $r = \|\vec{r}\|$ is:}{$n r^{n-1}\vec{r}$}{$n r^{n-2}\vec{r}$}{$r^n \vec{r}$}{$n r^n$}{1}

\mcqitem{18}{The value of $\nabla^2(1/r)$ for $r \neq 0$ is:}{$0$}{$1/r^2$}{$-1/r^3$}{$4\pi$}{1}

\mcqitem{19}{If $\vec{F} = \nabla \phi$, then the work done in moving a particle from point $A$ to point $B$ in the field is:}{$\phi(B) - \phi(A)$}{$\phi(A) - \phi(B)$}{$\phi(A)\phi(B)$}{$0$}{1}

\mcqitem{20}{The value of the integral $\int_0^{\pi/2} \int_0^{\pi/2} \sin(x+y) \, dx \, dy$ is:}{$1$}{$2$}{$0$}{$\pi$}{1}

\mcqitem{21}{Changing the order of integration of $\int_0^1 \int_0^x f(x,y) \, dy \, dx$ gives:}{$\int_0^1 \int_y^1 f(x,y) \, dx \, dy$}{$\int_0^1 \int_0^y f(x,y) \, dx \, dy$}{$\int_0^x \int_0^1 f(x,y) \, dx \, dy$}{$\int_y^1 \int_0^1 f(x,y) \, dx \, dy$}{1}

\mcqitem{22}{The total volume of a sphere of radius $a$ computed via triple integration is:}{$\frac{4}{3}\pi a^3$}{$4\pi a^2$}{$\frac{2}{3}\pi a^3$}{$\pi a^3$}{1}

\mcqitem{23}{The curl of the gradient of any scalar field $\phi$ is identically:}{$\vec{0}$}{$1$}{$\nabla^2 \phi$}{$\vec{r}$}{1}

\mcqitem{24}{The divergence of the curl of any vector field $\vec{F}$ is identically:}{$0$}{$1$}{$\nabla^2 \vec{F}$}{$\vec{F}$}{1}

\mcqitem{25}{A function $f(x,y)$ has a local minimum at $(a,b)$ if:}{$rt - s^2 > 0$ and $r > 0$}{$rt - s^2 > 0$ and $r < 0$}{$rt - s^2 < 0$}{$rt - s^2 = 0$}{1}

\mcqitem{26}{The unit normal vector to the surface $\phi(x,y,z) = c$ at a point is given by:}{$\frac{\nabla \phi}{\|\nabla \phi\|}$}{$\nabla \phi$}{$\frac{\vec{r}}{\|\vec{r}\|}$}{$\nabla \times \vec{r}$}{1}

\mcqitem{27}{The value of $\oint_C (x dy - y dx)$ over a closed curve enclosing area $A$ is:}{$A$}{$2A$}{$0$}{$-A$}{1}

\mcqitem{28}{If $u = x+y$ and $v = x-y$, the Jacobian $\frac{\partial(u,v)}{\partial(x,y)}$ is:}{$2$}{$-2$}{$0$}{$1$}{1}

\mcqitem{29}{The integral $\iint_R r \, dr \, d\theta$ represents the:}{Volume}{Area of region $R$}{Circumference}{Moment of Inertia}{1}

\mcqitem{30}{The value of $\iint_S \vec{r} \cdot \hat{n} \, dS$ over a closed surface bounding volume $V$ is:}{$V$}{$2V$}{$3V$}{$0$}{1}

\vspace{5mm}
\hrule
\vspace{5mm}

% ------------------------------------------------------------------------------
% PART B: 6 SUBJECTIVE QUESTIONS (ATTEMPT ANY 4)
% ------------------------------------------------------------------------------
\subsection*{PART B: Descriptive Section (Attempt Any 4 of 6) [40 Marks]}

\questionitem{31 (Unit 1)}{
    Find the eigenvalues and a complete set of linearly independent eigenvectors of the symmetric matrix:
    \[
    A = \begin{bmatrix} 6 & -2 & 2 \\ -2 & 3 & -1 \\ 2 & -1 & 3 \end{bmatrix}
    \]
    Verify that the eigenvectors corresponding to distinct eigenvalues are mutually orthogonal.
}{10}{CO1}{L3 Apply}

\questionitem{32 (Unit 2)}{
    Solve the third-order non-homogeneous ordinary differential equation by the operator method:
    \[
    (D^3 - 3D^2 + 4)y = e^{2x} + 40 \cos x
    \]
}{10}{CO2}{L3 Apply}

\questionitem{33 (Unit 3)}{
    \textbf{(a)} Test the convergence of the infinite positive series:
    \[
    \sum_{n=1}^\infty \frac{1 \cdot 3 \cdot 5 \cdots (2n-1)}{2 \cdot 4 \cdot 6 \cdots (2n)} \cdot \frac{1}{2n+1}
    \]
    \textbf{(b)} Determine whether the alternating series $\sum_{n=2}^\infty \frac{(-1)^n}{n (\ln n)^2}$ converges absolutely or conditionally.
}{10}{CO3}{L4 Analyze}

\questionitem{34 (Unit 4)}{
    A rectangular box open at the top is to have a volume of $32\text{ cm}^3$. Find the dimensions of the box that require the \textbf{minimum surface area} of material, using the \textbf{Method of Lagrange Multipliers}.
}{10}{CO4}{L3 Apply}

\questionitem{35 (Unit 5)}{
    \textbf{(a)} Change the order of integration and evaluate:
    \[
    \int_0^1 \int_{\sqrt{y}}^1 \sqrt{x^3 + 1} \, dx \, dy
    \]
    \textbf{(b)} Evaluate $\iiint_V z \, dV$ where $V$ is the solid upper hemisphere $x^2 + y^2 + z^2 \le a^2$ ($z \ge 0$).
}{10}{CO5}{L3 Apply}

\questionitem{36 (Unit 6)}{
    Verify the \textbf{Gauss Divergence Theorem} for the vector field:
    \[
    \vec{F} = 4x\hat{i} - 2y^2\hat{j} + z^2\hat{k}
    \]
    over the cylindrical solid bounded by $x^2 + y^2 = 4$, $z = 0$, and $z = 3$.
}{10}{CO6}{L4 Analyze}

\vspace{5mm}
\hrule
\vspace{5mm}

% ------------------------------------------------------------------------------
% COMPLETE STEP-BY-STEP SOLUTIONS
% ------------------------------------------------------------------------------
\subsection*{Solutions \& Marking Scheme (End-Term Set A)}

\subsubsection*{Part A: Answer Key \& Explanations}

\begin{center}
\begin{tabular}{|c|c||c|c||c|c||c|c||c|c||c|c|}
\hline
\textbf{Q} & \textbf{Ans} & \textbf{Q} & \textbf{Ans} & \textbf{Q} & \textbf{Ans} & \textbf{Q} & \textbf{Ans} & \textbf{Q} & \textbf{Ans} & \textbf{Q} & \textbf{Ans} \\
\hline
1 & \textbf{A} & 6 & \textbf{A} & 11 & \textbf{A} & 16 & \textbf{B} & 21 & \textbf{A} & 26 & \textbf{A} \\
2 & \textbf{A} & 7 & \textbf{B} & 12 & \textbf{A} & 17 & \textbf{B} & 22 & \textbf{A} & 27 & \textbf{B} \\
3 & \textbf{A} & 8 & \textbf{B} & 13 & \textbf{A} & 18 & \textbf{A} & 23 & \textbf{A} & 28 & \textbf{B} \\
4 & \textbf{A} & 9 & \textbf{D} & 14 & \textbf{B} & 19 & \textbf{A} & 24 & \textbf{A} & 29 & \textbf{B} \\
5 & \textbf{C} & 10 & \textbf{A} & 15 & \textbf{B} & 20 & \textbf{B} & 25 & \textbf{A} & 30 & \textbf{C} \\
\hline
\end{tabular}
\end{center}

\begin{solutionbox}[Explanations for Part A]
\textbf{Q.1 (Ans: A)} Let $x = r\cos\theta, y = r\sin\theta$: $\frac{xy}{\sqrt{x^2+y^2}} = \frac{r^2\cos\theta\sin\theta}{r} = r\cos\theta\sin\theta \to 0$ as $r\to 0$.\\[1.5mm]
\textbf{Q.4 (Ans: A)} $J = \begin{vmatrix} \cos\theta & -r\sin\theta \\ \sin\theta & r\cos\theta \end{vmatrix} = r(\cos^2\theta + \sin^2\theta) = r$.\\[1.5mm]
\textbf{Q.9 (Ans: D)} $\nabla \cdot \vec{r} = \frac{\partial x}{\partial x} + \frac{\partial y}{\partial y} + \frac{\partial z}{\partial z} = 1 + 1 + 1 = 3$.\\[1.5mm]
\textbf{Q.10 (Ans: A)} $\nabla \times \vec{r} = \vec{0}$ because the position vector field is purely irrotational.\\[1.5mm]
\textbf{Q.17 (Ans: B)} $\nabla(r^n) = n r^{n-1}\nabla r = n r^{n-1}\frac{\vec{r}}{r} = n r^{n-2}\vec{r}$.\\[1.5mm]
\textbf{Q.30 (Ans: C)} By Gauss Divergence Theorem, $\iint_S \vec{r} \cdot \hat{n} \, dS = \iiint_V (\nabla \cdot \vec{r}) \, dV = \iiint_V 3 \, dV = 3V$.
\end{solutionbox}

\subsubsection*{Part B: Detailed Subjective Solutions}

% Solution 31
\begin{solutionbox}[Solution to Question 31 (10 Marks)]
\textbf{Step 1: Characteristic Equation}
$A = \begin{bmatrix} 6 & -2 & 2 \\ -2 & 3 & -1 \\ 2 & -1 & 3 \end{bmatrix}$.
\begin{itemize}[leftmargin=6mm]
    \item $S_1 = \text{trace}(A) = 6 + 3 + 3 = 12$.
    \item $S_2 = M_{11} + M_{22} + M_{33} = (9 - 1) + (18 - 4) + (18 - 4) = 8 + 14 + 14 = 36$.
    \item $S_3 = \det(A) = 6(8) - (-2)(-6 + 2) + 2(2 - 6) = 48 + 2(-4) + 2(-4) = 48 - 8 - 8 = 32$.
\end{itemize}
Equation: $\lambda^3 - 12\lambda^2 + 36\lambda - 32 = 0$.
Factorization: $(\lambda - 2)(\lambda - 2)(\lambda - 8) = 0 \implies \mathbf{\lambda_1 = 8, \quad \lambda_2 = 2, \quad \lambda_3 = 2}$.

\textbf{Step 2: Eigenvector for $\lambda = 8$}
$(A - 8I)X = 0 \implies \begin{bmatrix} -2 & -2 & 2 \\ -2 & -5 & -1 \\ 2 & -1 & -5 \end{bmatrix} \begin{bmatrix} x_1 \\ x_2 \\ x_3 \end{bmatrix} = \begin{bmatrix} 0 \\ 0 \\ 0 \end{bmatrix}$.
From row 1: $-x_1 - x_2 + x_3 = 0$.
Using cross-multiplication on first two rows:
$\frac{x_1}{(-2)(-1) - (2)(-5)} = \frac{x_2}{(2)(-2) - (-2)(-1)} = \frac{x_3}{(-2)(-5) - (-2)(-2)} \implies \frac{x_1}{12} = \frac{x_2}{-6} = \frac{x_3}{6}$.
Dividing by 6 gives: $\mathbf{X_1 = \begin{bmatrix} 2 \\ -1 \\ 1 \end{bmatrix}}$.

\textbf{Step 3: Eigenvectors for Repeated Root $\lambda = 2$}
$(A - 2I)X = 0 \implies \begin{bmatrix} 4 & -2 & 2 \\ -2 & 1 & -1 \\ 2 & -1 & 1 \end{bmatrix} \begin{bmatrix} x_1 \\ x_2 \\ x_3 \end{bmatrix} = \begin{bmatrix} 0 \\ 0 \\ 0 \end{bmatrix}$.
All three rows are scalar multiples of $2x_1 - x_2 + x_3 = 0 \implies x_2 = 2x_1 + x_3$.
Here we have two free variables ($x_1$ and $x_3$):
\begin{itemize}[leftmargin=6mm]
    \item Let $x_1 = 1, x_3 = 0 \implies x_2 = 2$: $\mathbf{X_2 = \begin{bmatrix} 1 \\ 2 \\ 0 \end{bmatrix}}$.
    \item Let $x_1 = 0, x_3 = 1 \implies x_2 = 1$: $X_3' = \begin{bmatrix} 0 \\ 1 \\ 1 \end{bmatrix}$.
\end{itemize}
To ensure $X_3$ is orthogonal to $X_2$ (Gram-Schmidt):
Let $X_3 = \begin{bmatrix} a \\ b \\ c \end{bmatrix}$ satisfy $2a - b + c = 0$ and $X_2 \cdot X_3 = a + 2b = 0 \implies a = -2b$.
Then $2(-2b) - b + c = 0 \implies c = 5b$. Choosing $b = 1$: $\mathbf{X_3 = \begin{bmatrix} -2 \\ 1 \\ 5 \end{bmatrix}}$.

\textbf{Step 4: Orthogonality Verification}
$X_1 \cdot X_2 = (2)(1) + (-1)(2) + (1)(0) = 2 - 2 + 0 = 0$.
$X_1 \cdot X_3 = (2)(-2) + (-1)(1) + (1)(5) = -4 - 1 + 5 = 0$.
$X_2 \cdot X_3 = (1)(-2) + (2)(1) + (0)(5) = -2 + 2 + 0 = 0$.
All three eigenvectors are mutually orthogonal.

\begin{markingrubric}
\textbf{Mark Distribution:}
$\bullet$ Characteristic equation and roots $\lambda = 8, 2, 2$: \textbf{3 Marks}\\
$\bullet$ Eigenvector for $\lambda = 8$: \textbf{2 Marks}\\
$\bullet$ Linearly independent eigenvectors for $\lambda = 2$: \textbf{3 Marks}\\
$\bullet$ Proof of mutual orthogonality: \textbf{2 Marks}
\end{markingrubric}
\end{solutionbox}

% Solution 32
\begin{solutionbox}[Solution to Question 32 (10 Marks)]
\textbf{Step 1: Complementary Function}
Auxiliary equation: $m^3 - 3m^2 + 4 = 0$.
Test $m = -1$: $(-1)^3 - 3(-1)^2 + 4 = -1 - 3 + 4 = 0 \implies (m+1)$ is a factor.
$(m+1)(m^2 - 4m + 4) = 0 \implies (m+1)(m-2)^2 = 0 \implies m = -1, 2, 2$.
\[
\mathbf{y_c = c_1 e^{-x} + (c_2 + c_3 x)e^{2x}}
\]
\textbf{Step 2: Particular Integral for $e^{2x}$ ($y_{p_1}$)}
Since $m = 2$ is a root of multiplicity 2:
\[
y_{p_1} = \frac{1}{(D+1)(D-2)^2} e^{2x} = \frac{1}{(2+1)} \frac{1}{(D-2)^2} e^{2x} = \frac{1}{3} \frac{x^2}{2!} e^{2x} = \mathbf{\frac{1}{6} x^2 e^{2x}}
\]
\textbf{Step 3: Particular Integral for $40 \cos x$ ($y_{p_2}$)}
\[
y_{p_2} = \frac{1}{D^3 - 3D^2 + 4} (40 \cos x) = \frac{40}{D(D^2) - 3D^2 + 4} \cos x
\]
Replace $D^2$ with $-1^2 = -1$:
\[
= \frac{40}{D(-1) - 3(-1) + 4} \cos x = \frac{40}{-D + 7} \cos x = \frac{40}{7 - D} \cos x
\]
Multiply numerator and denominator by $(7 + D)$:
\[
= \frac{40(7 + D)}{49 - D^2} \cos x = \frac{40(7 + D)}{49 - (-1)} \cos x = \frac{40(7 + D)}{50} \cos x = \frac{4}{5}(7\cos x - \sin x)
\]
\textbf{Step 4: General Solution}
\[
\mathbf{y(x) = c_1 e^{-x} + (c_2 + c_3 x)e^{2x} + \frac{1}{6} x^2 e^{2x} + \frac{28}{5}\cos x - \frac{4}{5}\sin x}
\]

\begin{markingrubric}
\textbf{Mark Distribution:}
$\bullet$ Auxiliary equation and roots $m = -1, 2, 2$: \textbf{3 Marks}\\
$\bullet$ Particular integral $y_{p_1}$ with repeated root handling: \textbf{3 Marks}\\
$\bullet$ Particular integral $y_{p_2}$ for cosine: \textbf{3 Marks}\\
$\bullet$ Complete general solution: \textbf{1 Mark}
\end{markingrubric}
\end{solutionbox}

% Solution 33
\begin{solutionbox}[Solution to Question 33 (10 Marks)]
\textbf{Part (a): Convergence of the Positive Term Series (5 Marks)}
Let $u_n = \frac{1 \cdot 3 \cdot 5 \cdots (2n-1)}{2 \cdot 4 \cdot 6 \cdots (2n)} \cdot \frac{1}{2n+1}$.
Then:
\[
u_{n+1} = \frac{1 \cdot 3 \cdot 5 \cdots (2n-1)(2n+1)}{2 \cdot 4 \cdot 6 \cdots (2n)(2n+2)} \cdot \frac{1}{2n+3}
\]
Apply D'Alembert's Ratio Test:
\[
\frac{u_n}{u_{n+1}} = \frac{2n+2}{2n+1} \cdot \frac{2n+3}{2n+1} = \frac{4n^2 + 10n + 6}{4n^2 + 4n + 1}
\]
$\lim_{n\to\infty} \frac{u_n}{u_{n+1}} = 1$. The Ratio Test is \textbf{inconclusive}!
Apply \textbf{Raabe's Test}:
\[
L = \lim_{n\to\infty} n\left(\frac{u_n}{u_{n+1}} - 1\right) = \lim_{n\to\infty} n\left(\frac{4n^2 + 10n + 6 - 4n^2 - 4n - 1}{4n^2 + 4n + 1}\right)
\]
\[
= \lim_{n\to\infty} n\left(\frac{6n + 5}{4n^2 + 4n + 1}\right) = \lim_{n\to\infty} \frac{6n^2 + 5n}{4n^2 + 4n + 1} = \frac{6}{4} = \mathbf{\frac{3}{2}}
\]
Since $L = \frac{3}{2} > 1$, by Raabe's Test, the series is \textbf{strictly CONVERGENT}.

\textbf{Part (b): Alternating Series Convergence (5 Marks)}
Given $\sum_{n=2}^\infty (-1)^n a_n$ where $a_n = \frac{1}{n (\ln n)^2}$.
\begin{enumerate}[label=(\roman*)]
    \item Consider the series of absolute terms: $\sum_{n=2}^\infty \frac{1}{n (\ln n)^2}$.
    Apply Cauchy's Integral Test with $f(x) = \frac{1}{x (\ln x)^2}$:
    \[
    \int_2^\infty \frac{dx}{x (\ln x)^2} = \left[ -\frac{1}{\ln x} \right]_2^\infty = 0 - \left(-\frac{1}{\ln 2}\right) = \frac{1}{\ln 2} < \infty
    \]
    The integral converges $\implies$ The series of absolute values converges!
    \item Since the absolute series converges, the given alternating series is \textbf{ABSOLUTELY CONVERGENT}.
\end{enumerate}

\begin{markingrubric}
\textbf{Mark Distribution:}
$\bullet$ Part (a) Ratio test failure ($L=1$): \textbf{2 Marks}\\
$\bullet$ Part (a) Raabe's test evaluation showing $L = 3/2 > 1$ (Convergent): \textbf{3 Marks}\\
$\bullet$ Part (b) Integral test on absolute series: \textbf{3 Marks}\\
$\bullet$ Part (b) Conclusion of absolute convergence: \textbf{2 Marks}
\end{markingrubric}
\end{solutionbox}

% Solution 34
\begin{solutionbox}[Solution to Question 34 (10 Marks)]
\textbf{Step 1: Mathematical Formulation}
Let $x, y$ be the horizontal dimensions of the base and $z$ be the height.
The box is open at the top, so its surface area consists of the base and 4 sides:
\[
S(x,y,z) = xy + 2xz + 2yz
\]
Subject to the volume constraint:
\[
g(x,y,z) = xyz - 32 = 0
\]
\textbf{Step 2: Lagrange Multiplier Equations}
Set $\nabla S = \lambda \nabla g$:
\begin{align}
S_x &= \lambda g_x \implies y + 2z = \lambda yz \label{eq:lag1} \\
S_y &= \lambda g_y \implies x + 2z = \lambda xz \label{eq:lag2} \\
S_z &= \lambda g_z \implies 2x + 2y = \lambda xy \label{eq:lag3}
\end{align}
\textbf{Step 3: Solving the System}
Multiply \eqref{eq:lag1} by $x$ and \eqref{eq:lag2} by $y$:
\[
xy + 2xz = \lambda xyz
\]
\[
xy + 2yz = \lambda xyz
\]
Equating both gives:
\[
xy + 2xz = xy + 2yz \implies 2xz = 2yz \implies \mathbf{x = y} \quad (\text{since } z \neq 0)
\]
Now substitute $x = y$ into \eqref{eq:lag3}:
\[
2x + 2x = \lambda x^2 \implies 4x = \lambda x^2 \implies \lambda = \frac{4}{x}
\]
Substitute $\lambda = \frac{4}{x}$ into \eqref{eq:lag1} with $y = x$:
\[
x + 2z = \left(\frac{4}{x}\right) x z = 4z \implies x = 2z \implies \mathbf{z = \frac{x}{2}}
\]
\textbf{Step 4: Using Volume Constraint}
Substitute $y = x$ and $z = x/2$ into $xyz = 32$:
\[
(x)(x)\left(\frac{x}{2}\right) = 32 \implies \frac{x^3}{2} = 32 \implies x^3 = 64 \implies \mathbf{x = 4\text{ cm}}
\]
Hence:
\[
\mathbf{y = 4\text{ cm}}, \quad \mathbf{z = 2\text{ cm}}
\]
\textbf{Step 5: Minimum Surface Area}
\[
S_{\min} = (4)(4) + 2(4)(2) + 2(4)(2) = 16 + 16 + 16 = \mathbf{48\text{ cm}^2}
\]

\begin{markingrubric}
\textbf{Mark Distribution:}
$\bullet$ Objective function $S$ and constraint $g$: \textbf{2 Marks}\\
$\bullet$ Lagrange system setup $\nabla S = \lambda \nabla g$: \textbf{3 Marks}\\
$\bullet$ Solving relation $x = y = 2z$: \textbf{3 Marks}\\
$\bullet$ Final dimensions $4 \times 4 \times 2\text{ cm}$ and $S = 48\text{ cm}^2$: \textbf{2 Marks}
\end{markingrubric}
\end{solutionbox}

% Solution 35
\begin{solutionbox}[Solution to Question 35 (10 Marks)]
\textbf{Part (a): Change of Order of Integration (5 Marks)}
Given: $I = \int_0^1 \int_{\sqrt{y}}^1 \sqrt{x^3 + 1} \, dx \, dy$.
Limits: $y \in [0, 1]$, $x \in [\sqrt{y}, 1] \implies y = x^2$.
In reversed order (vertical strips):
$x \in [0, 1]$, $y \in [0, x^2]$.
\[
I = \int_0^1 \int_0^{x^2} \sqrt{x^3 + 1} \, dy \, dx = \int_0^1 \sqrt{x^3 + 1} [y]_0^{x^2} \, dx = \int_0^1 x^2 \sqrt{x^3 + 1} \, dx
\]
Let $u = x^3 + 1 \implies du = 3x^2 \, dx \implies x^2 \, dx = \frac{du}{3}$.
When $x = 0, u = 1$; when $x = 1, u = 2$:
\[
I = \frac{1}{3} \int_1^2 u^{1/2} \, du = \frac{1}{3} \left[ \frac{u^{3/2}}{3/2} \right]_1^2 = \frac{2}{9} (2^{3/2} - 1) = \mathbf{\frac{2}{9}(2\sqrt{2} - 1)}
\]

\textbf{Part (b): Triple Integral in Spherical Coordinates (5 Marks)}
Region $V$: Upper hemisphere of radius $a$ ($z \ge 0$).
In spherical polar coordinates: $z = \rho \cos\phi$, $dV = \rho^2 \sin\phi \, d\rho \, d\phi \, d\theta$.
Limits: $\rho \in [0, a]$, $\phi \in [0, \pi/2]$ (upper half), $\theta \in [0, 2\pi]$.
\[
\iiint_V z \, dV = \int_0^{2\pi} d\theta \cdot \int_0^{\pi/2} \cos\phi \sin\phi \, d\phi \cdot \int_0^a \rho^3 \, d\rho
\]
\begin{itemize}[leftmargin=6mm]
    \item $\int_0^{2\pi} d\theta = 2\pi$
    \item $\int_0^{\pi/2} \sin\phi \cos\phi \, d\phi = \left[ \frac{\sin^2\phi}{2} \right]_0^{\pi/2} = \frac{1}{2}$
    \item $\int_0^a \rho^3 \, d\rho = \frac{a^4}{4}$
\end{itemize}
Multiplying:
\[
\iiint_V z \, dV = 2\pi \left(\frac{1}{2}\right) \left(\frac{a^4}{4}\right) = \mathbf{\frac{\pi a^4}{4}}
\]

\begin{markingrubric}
\textbf{Mark Distribution:}
$\bullet$ Part (a) reversed limits and integration yielding $\frac{2}{9}(2\sqrt{2}-1)$: \textbf{5 Marks}\\
$\bullet$ Part (b) spherical bounds and integration yielding $\frac{\pi a^4}{4}$: \textbf{5 Marks}
\end{markingrubric}
\end{solutionbox}

% Solution 36
\begin{solutionbox}[Solution to Question 36 (10 Marks)]
\textbf{Step 1: Statement of Gauss Divergence Theorem}
\[
\iint_S \vec{F} \cdot \hat{n} \, dS = \iiint_V (\nabla \cdot \vec{F}) \, dV
\]
Given $\vec{F} = 4x\hat{i} - 2y^2\hat{j} + z^2\hat{k}$ over cylinder $x^2 + y^2 \le 4, 0 \le z \le 3$.

\textbf{Step 2: Volume Integral Evaluation}
Compute $\nabla \cdot \vec{F}$:
\[
\nabla \cdot \vec{F} = \frac{\partial}{\partial x}(4x) + \frac{\partial}{\partial y}(-2y^2) + \frac{\partial}{\partial z}(z^2) = 4 - 4y + 2z
\]
Transform to cylindrical coordinates: $x = r\cos\theta, y = r\sin\theta, z = z$, $dV = r \, dz \, dr \, d\theta$.
Bounds: $r \in [0, 2]$, $\theta \in [0, 2\pi]$, $z \in [0, 3]$.
\[
\iiint_V (4 - 4r\sin\theta + 2z) r \, dz \, dr \, d\theta
\]
Due to periodicity, $\int_0^{2\pi} \sin\theta \, d\theta = 0$, so the middle term vanishes!
\[
= \int_0^{2\pi} d\theta \int_0^2 r \, dr \int_0^3 (4 + 2z) dz = (2\pi) \left[ \frac{r^2}{2} \right]_0^2 \left[ 4z + z^2 \right]_0^3
\]
\[
= (2\pi)(2)(12 + 9) = 4\pi (21) = \mathbf{84\pi}
\]

\textbf{Step 3: Surface Integral Evaluation (3 Surfaces)}
\begin{enumerate}[label=(\roman*)]
    \item \textbf{Bottom Surface $S_1$ ($z = 0$):} $\hat{n} = -\hat{k}, z = 0 \implies \vec{F} \cdot \hat{n} = -z^2 = 0 \implies \iint_{S_1} = 0$.
    \item \textbf{Top Surface $S_2$ ($z = 3$):} $\hat{n} = \hat{k}, z = 3 \implies \vec{F} \cdot \hat{n} = z^2 = 9$.
    \[
    \iint_{S_2} 9 \, dS = 9 \times \text{Area}(\text{circle of radius 2}) = 9(\pi \cdot 2^2) = 36\pi
    \]
    \item \textbf{Curved Surface $S_3$ ($r = 2$):}
    Normal $\hat{n} = \frac{x\hat{i} + y\hat{j}}{2}$.
    \[
    \vec{F} \cdot \hat{n} = \frac{4x^2 - 2y^3}{2} = 2x^2 - y^3 = 2(2\cos\theta)^2 - (2\sin\theta)^3 = 8\cos^2\theta - 8\sin^3\theta
    \]
    $dS = 2 \, d\theta \, dz$:
    \[
    \iint_{S_3} \vec{F} \cdot \hat{n} \, dS = \int_0^3 dz \int_0^{2\pi} (8\cos^2\theta - 8\sin^3\theta) 2 \, d\theta
    \]
    Since $\int_0^{2\pi} \sin^3\theta d\theta = 0$ and $\int_0^{2\pi} \cos^2\theta d\theta = \pi$:
    \[
    = (3)(2)(8\pi) = 48\pi
    \]
\end{enumerate}
Total surface integral:
\[
\iint_S \vec{F} \cdot \hat{n} \, dS = 0 + 36\pi + 48\pi = \mathbf{84\pi}
\]
Since Volume Integral ($84\pi$) = Surface Integral ($84\pi$), **Gauss Divergence Theorem is verified!**

\begin{markingrubric}
\textbf{Mark Distribution:}
$\bullet$ Stating theorem and computing $\nabla \cdot \vec{F} = 4 - 4y + 2z$: \textbf{2 Marks}\\
$\bullet$ Evaluating volume integral to obtain $84\pi$: \textbf{3 Marks}\\
$\bullet$ Evaluating surface integrals across top, bottom, and curved faces: \textbf{4 Marks}\\
$\bullet$ Equating and concluding verification: \textbf{1 Mark}
\end{markingrubric}
\end{solutionbox}

\newpage
